\section{Convergence, contour comparison, and residues}\label{sec:convergence}

\subsection{The vertical-line estimate}

Let $s=\sigma+\ii t$ and let $\phi=\Arg z$ on the certified branch.  Applying
the two Gamma bounds \eqref{eq:gamma-upper}--\eqref{eq:recip-gamma-upper} to
each factor of \eqref{eq:phi-minus} gives a computable constant $C>0$ such
that
\begin{equation}\label{eq:master-vertical-bound}
 \abs{\Phi(\sigma+\ii t)z^{\sigma+\ii t}}
 \le C\abs z^\sigma(1+\abs t)^A
      \exp\!\bigl(-[\pi\delta-\abs\phi]\abs t\bigr),
\end{equation}
where
\begin{equation}\label{eq:A-exponent}
 A=\Re\nu+(p-q)\left(\sigma+\frac12\right).
\end{equation}
Indeed, each numerator Gamma contributes $\e^{-\pi\abs t/2}$ and each
reciprocal denominator Gamma contributes $\e^{\pi\abs t/2}$, leaving
$\e^{-\pi\delta\abs t}$.  The powers of $1+\abs t$ from all $a$-parameters and
$b$-parameters sum exactly to \eqref{eq:A-exponent}.  Finally,
$\abs{z^s}=\abs z^\sigma\e^{-t\phi}$, whose symmetric upper bound is
$\abs z^\sigma\e^{\abs\phi\abs t}$.

\begin{theorem}[Constructive vertical convergence]\label{thm:vertical-convergence}
Suppose the line $\Re s=\sigma$ has certified pole clearance.
\begin{enumerate}[label=(\alph*)]
\item If $\epsilon:=\pi\delta-\abs{\Arg z}>0$ with a rational lower bound,
      then the vertical integral converges with a primitive-recursive tail
      modulus.
\item On the boundary $\epsilon=0$, it converges whenever a rational
      $\eta>0$ certifies $A\le-1-\eta$; equivalently,
      $(q-p)(\sigma+1/2)\ge\Re\nu+1+\eta$.
\end{enumerate}
These are the quantitative versions of the first contour conditions on the
source page.
\end{theorem}
\begin{proof}
In case (a), put $M=\max(0,\lceil A\rceil)$.  For
$T\ge(4M/\epsilon)^2+1$, the elementary estimate
$\log(1+t)\le\sqrt t$ gives
$(1+t)^M\le\e^{\epsilon t/2}$.  Therefore the two tails are at most
\[
 \frac{4C\abs z^\sigma}{\epsilon}\e^{-\epsilon T/2}.
\]
A bounded search using the finite Taylor enclosure for $\exp$ finds a $T(k)$
for which this is below $2^{-k}$.  In case (b),
\[
 \int_T^\infty C(1+t)^{-1-\eta}\dd t
 \le \frac C\eta T^{-\eta}.
\]
Writing $\eta=r/s$ rational, choose $T$ by the integer inequality
$T^r\ge(2^kC/\eta)^s$.  Integer powering and bounded root search are primitive
recursive.  The finite central segment is covered by rational quadrature.
\end{proof}

\subsection{Right residues and Slater's formula}

Assume no two $b_1,\ldots,b_m$ differ by an integer.  The right poles are
$s=b_h+k$, $h=1,\ldots,m$, $k\ge0$.  The residue of
$\Gamma(b_h-s)$ with respect to $s$ is
$(-1)^{k+1}/k!$.  The right loop is negatively oriented, so its extra minus
sign changes this to $(-1)^k/k!$.  Apply
\begin{equation}\label{eq:gamma-shift-identities}
 \Gamma(x+k)=\Gamma(x)\poch{x}{k},
 \qquad
 \Gamma(x-k)=\frac{(-1)^k\Gamma(x)}{\poch{1-x}{k}}
\end{equation}
to every remaining factor.

\begin{theorem}[Right-residue expansion / Slater]\label{thm:slater-right}
If $p<q$, or if $p=q$ and $\abs z<1$ with rational margin, and the right poles
are nonconfluent, then
\begin{align}
&G_{p,q}^{m,n}\!\left(\left.\begin{matrix}\mathbf a\\\mathbf b\end{matrix}\right|z\right)
 =\sum_{h=1}^{m}C_h z^{b_h}
 {}_pF_{q-1}\!\left(
 \left.\begin{matrix}1+b_h-\mathbf a\\
 (1+b_h-\mathbf b)^{\!*}
 \end{matrix}\right|(-1)^{p-m-n}z\right),\label{eq:slater-right}\\
&C_h=
 \frac{\prod_{j=1,\,j\ne h}^{m}\Gamma(b_j-b_h)
       \prod_{j=1}^{n}\Gamma(1+b_h-a_j)}
      {\prod_{j=m+1}^{q}\Gamma(1+b_h-b_j)
       \prod_{j=n+1}^{p}\Gamma(a_j-b_h)}.
\end{align}
The star means that the lower parameter with index $h$ is omitted.  The equality
has a primitive-recursive modulus on every certified compact branch domain.
\end{theorem}
\begin{proof}
The residue calculation above gives the $k$th term.  Substituting
\eqref{eq:gamma-shift-identities} and collecting the signs yields precisely the
coefficient recurrence \eqref{eq:hypergeom}.  When $p<q$, the series on the
right has at least as many lower parameters as upper parameters and converges
for all bounded $z$ by \cref{thm:hyp-conv}; when $p=q$, it has radius one and
the certified margin supplies a geometric tail.  A finite right-closing
rectangle contains finitely many poles.  The residue theorem for that polygon
is a finite triangulation identity plus local Laurent-jet calculations.  Gamma
ratio bounds make the outer side smaller than the same hypergeometric tail.
Given $k$, choose the rectangle and series cutoffs so that both errors are below
$2^{-k-2}$; the finite residue sums agree exactly.  Ball transitivity proves
\eqref{eq:slater-right}.
\end{proof}

\begin{corollary}[Empty right family]\label{cor:mzero}
For the right-loop definition,
$G_{p,q}^{0,n}=0$ when $p<q$, or when $p=q$ and $\abs z<1$.
\end{corollary}
\begin{proof}
Every finite right contour contains no poles, so its integral is zero by the
finite Cauchy theorem.  The closing-side error tends below every requested
dyadic tolerance.  Equivalently, the sum in \eqref{eq:slater-right} is empty.
\end{proof}

\subsection{Left residues}

The poles enclosed by the left loop are $s=a_h-1-k$ for $h\le n$.  The same
calculation, or inversion from \cref{thm:inversion}, gives the reciprocal
formula.

\begin{theorem}[Left-residue expansion]\label{thm:slater-left}
If $p>q$, or if $p=q$ and $\abs z>1$ with rational margin, and the left poles
are nonconfluent, then
\begin{align}
&G_{p,q}^{m,n}\!\left(\left.\begin{matrix}\mathbf a\\\mathbf b\end{matrix}\right|z\right)
 =\sum_{h=1}^{n}D_h z^{a_h-1}
 {}_qF_{p-1}\!\left(
 \left.\begin{matrix}1-a_h+\mathbf b\\
 (1-a_h+\mathbf a)^{\!*}
 \end{matrix}\right|(-1)^{q-m-n}z^{-1}\right),\label{eq:slater-left}\\
&D_h=
 \frac{\prod_{j=1,\,j\ne h}^{n}\Gamma(a_h-a_j)
       \prod_{j=1}^{m}\Gamma(1-a_h+b_j)}
      {\prod_{j=n+1}^{p}\Gamma(1-a_h+a_j)
       \prod_{j=m+1}^{q}\Gamma(a_h-b_j)}.
\end{align}
\end{theorem}
The proof is the reflected version of \cref{thm:slater-right}.  The argument of
the hypergeometric series has norm below one when $p=q$ because it is
proportional to $z^{-1}$.

\begin{corollary}[Empty left family]\label{cor:nzero}
For the left-loop definition, $G_{p,q}^{m,0}=0$ when $p>q$, or when $p=q$ and
$\abs z>1$.
\end{corollary}

\begin{theorem}[Balanced unit-circle boundary]\label{thm:balanced-unit-circle}
Suppose $p=q$, $\abs z=1$, and a rational margin $\eta>0$ certifies
$\Re\nu\le-1-\eta$.  Then both the right-loop and left-loop residue
presentations converge whenever their pole-separation data are admissible.
For every fixed confluent multiplicity they have a primitive-recursive
$p$-series tail modulus as well.
\end{theorem}
\begin{proof}
Consider one right Slater summand.  Its $k$th coefficient, apart from its fixed
prefactor and the unit-modulus phase, is a quotient of Pochhammer symbols.
Writing each Pochhammer symbol as a Gamma quotient and using the fixed-strip
Stirling bounds gives computable $C,N$ such that
\begin{equation}\label{eq:unit-circle-coeff}
 \abs{t_{h,k}}\le C(k+1)^{\Re\nu+\eta/2}
 \le C(k+1)^{-1-\eta/2}\qquad(k\ge N).
\end{equation}
The exponent follows by direct cancellation:
\[
 \sum_{i=1}^{p}\Re(1+b_h-a_i)
 -\sum_{j\ne h}\Re(1+b_h-b_j)-1=\Re\nu.
\]
The Stirling statement can alternatively be checked directly from the squared
term-ratio polynomial: its leading coefficient is $1$ and its next coefficient
is $2\Re\nu$, so bounded search finds the same $N$ with the prescribed margin.
For $a=\eta/2$,
\[
 \sum_{k\ge N}(k+1)^{-1-a}\le a^{-1}N^{-a};
\]
when $a=r/s$ is rational, the required cutoff is obtained from an integer
inequality after raising to the power $s$.  This is a primitive-recursive
$p$-series certificate.

A pole group of fixed multiplicity contributes an additional fixed polynomial
in $\Log z$, harmonic sums, and $\log(k+1)$.  Since bounded search proves
$(1+\log(k+1))^d\le(k+1)^{\eta/4}$ eventually for every fixed $d$, the same
argument leaves exponent at most $-1-\eta/4$.  The left expansion is identical
after inversion, because $\abs{z^{-1}}=1$ and the same $\nu$ occurs.  The
closing-side estimates are bounded by these residue tails, completing the two
loop integrals.
\end{proof}

\subsection{Confluent poles}

If $b_j-b_h\in\Z$, several right-pole sequences overlap.  The simple constants
$C_h$ then contain divergent Gamma factors, but the original contour integral
need not diverge.  Group the overlapping poles and apply the jet algorithm
\cref{def:jet}.  A multiplicity-$r$ group contributes
\begin{equation}\label{eq:power-log-normal}
 z^{\beta+k}\sum_{j=0}^{r-1}c_{k,j}(\Log z)^j.
\end{equation}
The coefficients $c_{k,j}$ are finite expressions in Gamma jets, factorials,
and harmonic sums.  Their ratios are bounded by a fixed polynomial in $k$
times the nonconfluent majorant, so the same convergence domains hold.  The
left case is identical.  Thus the ``no confluent poles'' hypothesis belongs to
the short closed form \eqref{eq:slater-right}, not to the existence of the
$G$-function.

\subsection{Agreement of contour types and analyticity}

\begin{theorem}[Contour agreement]\label{thm:contour-agreement}
Whenever two of the vertical, right-loop, and left-loop presentations satisfy
their quantitative hypotheses on the same branch domain, they are ball-equal
to every precision.
\end{theorem}
\begin{proof}
Truncate both contours and join them by a rational rectangle avoiding every
pole.  Finite triangulation cancels interior edges.  The remaining difference
is the sum of residues in the strip plus the outer sides.  If no poles lie
between the contours, the residue sum is empty.  Otherwise both presentations
are compared with the same finite residue normal form.  The vertical estimate
\eqref{eq:master-vertical-bound} and the right/left series tails make all outer
pieces arbitrarily small with primitive-recursive cutoffs.
\end{proof}

\begin{theorem}[Quantitative analyticity]\label{thm:analyticity}
On a compact rational domain that avoids $z=0$, the chosen branch cut, and in
the balanced case the singular point selected by the differential equation,
the certified $G$-name has Taylor coefficients and a primitive-recursive
uniform remainder modulus.  Every externally fixed derivative order is
therefore represented by a primitive-recursive name.
\end{theorem}
\begin{proof}
In a residue domain, differentiate the uniformly convergent power-log series;
polynomial factors in the summation index preserve the ratio tail.  In a
vertical domain, differentiate $z^s$ an externally fixed number of times and
use \eqref{eq:master-vertical-bound} with a larger polynomial exponent.  A
finite Cauchy estimate on a rational circle converts derivative bounds into a
Taylor remainder.  All precision choices are bounded searches.
\end{proof}
